\subsection{Feature formats and pipeline}
We compare three feature formats of increasing cost, one per model line, so that Section~\ref{sec:formats} can measure the trade-off directly:
\begin{enumerate}
  \item \textbf{Sparse bag-of-n-grams with TF-IDF weighting}, used by linear classifiers;
  \item \textbf{static pre-trained word embeddings} (300-dimensional GloVe\cite{pennington2014glove}), used by an averaged-embedding classifier and by a BiLSTM;
  \item \textbf{subword tokens fed to a pre-trained Transformer} (RoBERTa-base), whose contextual representation is fine-tuned end-to-end. This is the format of our final model.
\end{enumerate}
Figure~\ref{fig:b-pipeline} summarizes the pipeline. All three formats share the loading, auditing, splitting and cleaning steps described below.

\begin{figure}[tbp]
\centering
\begin{tikzpicture}[
  node distance=4mm and 6mm,
  box/.style={draw, rounded corners=2pt, align=center, font=\footnotesize, minimum height=8mm, text width=27mm, fill=gray!8},
  fbox/.style={box, text width=36mm, fill=gray!20},
  arr/.style={-{Stealth[length=2mm]}, thick}]
\node[box] (load) {Load \texttt{train.json}\\(7{,}401) / \texttt{test.json} (1{,}851)};
\node[box, right=of load] (audit) {Audit: imbalance,\\lengths, duplicates};
\node[box, right=of audit] (split) {Grouped stratified split\\5{,}909 train / 1{,}492 val};
\node[box, right=of split] (clean) {Light cleaning\\(HTML, whitespace)};
\node[fbox, below=9mm of audit, xshift=-6mm] (tfidf) {\textbf{TF-IDF}: lowercase, word 1--2-grams (+ char 2--5-grams), sublinear tf, L2 $\rightarrow$ sparse $x\in\mathbb{R}^{50\text{k}}$};
\node[fbox, right=of tfidf] (glove) {\textbf{GloVe}: lowercase tokens $\rightarrow$ vocabulary ids $\rightarrow$ 300-d vectors (mean or sequence, $L\le256$)};
\node[fbox, right=of glove] (bpe) {\textbf{RoBERTa}: cased byte-level BPE $\rightarrow$ ids + mask, $L\le256$};
\draw[arr] (load) -- (audit);
\draw[arr] (audit) -- (split);
\draw[arr] (split) -- (clean);
\draw[arr] (clean.south) |- ++(0,-4mm) -| (tfidf.north);
\draw[arr] (clean.south) |- ++(0,-4mm) -| (glove.north);
\draw[arr] (clean.south) -- ++(0,-4mm) -| (bpe.north);
\end{tikzpicture}
\caption{Data preprocessing pipeline. Every vectorizer, vocabulary and model is fitted on the training part only; the validation split is used only for model selection and the test set only for the final predictions.}
\label{fig:b-pipeline}
\end{figure}

\subsection{Data audit}
Both files are read as JSON tables.
The training set has 7{,}401 reviews with a binary label (6{,}319 positive, 1{,}082 negative). The test set has 1{,}851 unlabeled reviews. No review is missing or empty.
Review length varies widely: the median is 42 words, 75\% of reviews have at most 70 words, 97\% at most 178, and the longest has 1{,}066 words.
We also checked for duplicates, comparing texts after lower-casing and collapsing whitespace:
\begin{itemize}
  \item 154 distinct texts occur more than once in the training set, covering 373 rows, and every copy of a text has the same label;
  \item 101 test reviews are identical to a training review;
  \item 21 test reviews duplicate another test review.
\end{itemize}

\subsection{Grouped, stratified validation split}
Every model is tuned on one shared validation split.
If duplicates were split at random, copies of the same review would land on both sides, and the validation score would partly measure memorization.
We therefore group rows by their normalized text and take fold~0 of a 5-fold \texttt{StratifiedGroupKFold} (seed 42).
This yields 5{,}909 training and 1{,}492 validation reviews. The validation set contains 218 negatives (14.6\%, the same ratio as the full set), and no text appears on both sides.
Cross-validation inside the training part, including the internal folds used to calibrate the SVM, uses the same grouping.
The test set is used only to produce \texttt{submission.csv}.

\subsection{Text cleaning}
We unescape HTML entities, collapse runs of whitespace and strip leading and trailing spaces.
We deliberately do \emph{not}:
\begin{itemize}
  \item remove stop words, because standard stop lists contain ``not'', ``no'' and ``but'', which carry sentiment;
  \item stem, or strip punctuation such as ``!'' and ``?''.
\end{itemize}
Text is lower-cased for TF-IDF and GloVe, and kept in its original case for RoBERTa, whose tokenizer is case-sensitive.

\subsection{TF-IDF features (L1)}
The text is tokenized with a pattern that keeps single-character tokens (e.g.\ the ``t'' of ``don't'') as well as ``!'' and ``?''.
We extract unigrams and bigrams; bigrams capture short negations such as ``not good'' and ``not worth''.
Each term $t$ in review $d$ is weighted with sublinear term frequency and smoothed inverse document frequency,
\begin{equation}
  w_{t,d} = \bigl(1 + \log \mathrm{tf}_{t,d}\bigr)\cdot\Bigl(\log\frac{1+N}{1+\mathrm{df}_t} + 1\Bigr),
\end{equation}
where $N$ is the number of training reviews and $\mathrm{df}_t$ the number of reviews containing $t$. Each review vector is then L2-normalized.
The vectorizer is fitted on the training rows only.
Three settings were tuned by grouped 5-fold cross-validation:
\begin{itemize}
  \item the minimum document frequency (1, 2 or 3);
  \item the vocabulary cap (10k, 30k, 50k or unlimited);
  \item whether to append character 2--5-grams, which can absorb misspellings.
\end{itemize}
The selected representation has 50{,}000 word 1--2-gram features. It is very sparse: 86 non-zeros per review on average, a density of 0.17\%.

\subsection{Static word embeddings (L2)}
The L2 tokenizer lowercases and unescapes the text, retains contractions and
the punctuation tokens \texttt{!} and \texttt{?}, and keeps up to 256 tokens.
The vocabulary is built only on the training portion, with minimum frequency
two and special \texttt{<pad>} and \texttt{<unk>} entries. Of 6,796 ordinary
training vocabulary entries, 6,359 have a vector in the 300-dimensional
\texttt{glove.6B} table\cite{pennington2014glove}.
On validation, 3.17\% of token occurrences map to \texttt{<unk>}, and
4.95\% have no usable pretrained vector after applying the vocabulary.
Mean GloVe averages only the covered vectors. BiLSTM gives each in-vocabulary
entry missing a GloVe vector and the shared unknown entry a random initial
vector. Padding is excluded from both directions and the pooled statistics.

\subsection{Subword tokens for a pre-trained Transformer (L3)}
After cleaning, reviews retain their original casing. RoBERTa's
byte-level BPE tokenizer converts the text to token IDs and adds the
\texttt{<s>} and \texttt{</s>} tokens. The 256-token limit includes
these special tokens. When a review is too long, the tokenizer keeps
the beginning of the text. Each batch is padded to the length of its
longest sequence, and the attention mask distinguishes input tokens
from padding. Section~\ref{sec:models} describes how the model uses the first token's
representation for classification.

\subsection{Labels and class imbalance}
Labels are used as given (0 = negative, 1 = positive).
Each model outputs $P(y{=}1\mid x)$. We report metrics at threshold~0.5 and at thresholds tuned on the validation split: one maximizing accuracy ($\tau_{\mathrm{acc}}$, used for the submission) and one maximizing Macro-F1 ($\tau_{\mathrm{F1}}$).
The RoBERTa runs used unweighted cross-entropy. Section~\ref{sec:hparams} describes the parameter settings evaluated for the other models.
